\documentclass[11pt]{article}
\usepackage[a4paper,margin=26mm]{geometry}
\usepackage{amsmath,amssymb,amsthm,booktabs,hyperref}
\newtheorem{theorem}{Theorem}
\newtheorem{lemma}[theorem]{Lemma}
\title{The primitive rank-two geometries of $M_{11}$}
\author{}\date{}
\begin{document}\maketitle
This note treats Tables 3 and 4 of Leemans and Rodrigues,
\emph{On primitive geometries of rank two}, Art Discrete Appl. Math.
2 (2019), P1.08. The fixed journal edition has twelve pages. Its numerical
suborbit table and its total of $37$ unordered subgroup-pair orbits were
already recovered by the common character and mark calculation. Here the
remaining task is to recover the individual graph invariants and design
parameters from concrete subgroup pairs. The earlier calculation is an
independent comparison, not an input prescribing the number of pairs.

\section{The finite input and its scope}
We use the acquired degree-eleven model
\[
 G=\langle(1\,2\,\cdots\,11),\ (3\,7\,11\,8)(4\,10\,5\,6)\rangle
\]
and its retained array of $7920$ elements. The accepted numerical subgroup
census and structural proofs used for L054 supply $39$ subgroup-class
representatives. Completeness of that imported census is an explicit
premise. Its five maximal classes, in their retained order, are
\[
\begin{array}{c|rrrrr}
\text{class}&38&37&35&34&27\\
\text{order}&720&660&144&120&48\\
\text{index}&11&12&55&66&165.
\end{array}
\]
The names $M_{10}$, $L_2(11)$, $M_9:2$, $S_5$ and $Q_8:S_3$ are imported
from the actual accepted structural proofs for these literal representatives.
The last name is identified with $\mathrm{GL}_2(3)$ below. We also retain
the classical fact $\operatorname{Out}(M_{11})=1$. Thus inner conjugacy
of unordered pairs is exactly the equivalence defined by the source.
No restriction such as the intersection property $(\mathrm{IP})_2$ is imposed.

\begin{lemma}
Let $G$ be a finite group with a complete list of maximal subgroup classes,
and let $\Omega$ be the union of their conjugacy orbits. If a supplied
partition of $\binom{\Omega}{2}$ is closed under generators of $G$, and
each part is connected under those generator actions, its parts are exactly
the conjugacy orbits of unordered pairs of distinct maximal subgroups.
\end{lemma}
\begin{proof}
Closure under a generator implies closure under its inverse, since the
part is finite and the generator acts injectively. Thus each part is
$G$-invariant and is a union of orbits. Connectivity means that every member
can be reached from the chosen representative by generators and inverses,
so that the part lies in one orbit. The two inclusions give equality.
\end{proof}
The certificate retains every maximal subgroup as an actual subset of $G$,
every conjugator from its class representative, the two generator actions,
and the complete partition of the unordered-pair domain. Consequently
completeness is checked on a finite domain, rather than inferred from
agreement with a printed total.

\section{Cosets, distances and cycles}
For a retained pair $H,K$, construct its full right-coset partitions.
The edge set is
\[
                 \{(Hg,Kg):g\in G\}.
\]
These are exactly the pairs of cosets with nonempty intersection: an element
in the intersection supplies $g$, and conversely the displayed two cosets
both contain $g$. An edge is represented by exactly $|H\cap K|$ elements,
so there are $|G|/|H\cap K|$ edges. Right multiplication is transitive on
each vertex type and on the edges. Our permutations use this right-action
convention, exactly as the source does. If a prior component uses left
cosets, inversion supplies the precise bridge
$Hg\mapsto g^{-1}H$: two right cosets intersect if and only if their inverse
left cosets intersect. It preserves types and distances, while right
multiplication by $a$ becomes left multiplication by $a^{-1}$. No exchange
of the two types is part of this bridge.

Distinct maximal subgroups generate $G$. The incidence graph is therefore
connected. Its two degrees are $[H:H\cap K]$ and $[K:H\cap K]$, both at
least two, so it is firm. The action on each type is primitive by maximality.
Thus every retained pair gives a geometry in the source's class.
All these assertions also have literal finite checks in the certificate.

The acquisition computes all vertex eccentricities and the global girth.
The independent verifier uses one vertex of each type for the eccentricity
and a different shortest-cycle calculation. For a chosen root $r$, delete
$r$, initialize a separate search region at each neighbor of $r$, and run
a simultaneous breadth-first search. A meeting of two regions gives a
shortest path between two different neighbors. Its length plus two is the
length of a shortest cycle through $r$. Conversely every cycle through $r$
gives such a path after deleting $r$. Each type is transitive under the
explicit right action of $G$, and any shortest cycle contains both types.
Hence each chosen root lies on a cycle of global minimum length. This
justifies the rooted verification; a single arbitrary rooted breadth-first
collision would not, without further argument, certify global girth.
The gonality is half the girth, as in the source.

For the intersection name, the certificate gives a literal $g\in G$ and a
retained subgroup class $L$ such that $H\cap K=g^{-1}Lg$. Equality is checked
on complete element sets. Only then is the accepted structural proof for
$L$ imported. Equal subgroup orders are never used as a structural test.
When the source prints an integer alone, the recovered assertion is only
that intersection order.

\section{The order-forty-eight name}
The acquired representative of class $27$ is generated by
\[
 a=(1\,2\,9\,4\,6\,5\,3\,10)(7\,11),\qquad
 b=(2\,3)(4\,10)(5\,9)(8\,11).
\]
A complete finite isomorphism sends these generators to
\[
 A=\begin{pmatrix}2&1\\1&0\end{pmatrix},\qquad
 B=\begin{pmatrix}1&2\\0&2\end{pmatrix}
       \quad\text{over }\mathbf F_3.
\]
The certificate lists one matrix image for each of the $48$ actual
permutations. The verifier checks that the images are precisely all
invertible $2$ by $2$ matrices, that identity maps to identity, and that
$f(xy)=f(x)f(y)$ for all $48^2=2304$ pairs. This is a bijective homomorphism,
and proves the literal $\mathrm{GL}_2(3)$ identification.

For comparison with the accepted extension structure, put
\[
 i=\begin{pmatrix}0&1\\2&0\end{pmatrix},\quad
 j=\begin{pmatrix}1&1\\1&2\end{pmatrix},\quad
 r=\begin{pmatrix}1&1\\0&1\end{pmatrix},\quad
 s=\begin{pmatrix}1&0\\0&2\end{pmatrix}.
\]
Here $\langle i,j\rangle\cong Q_8$, with $i^2=j^2=-I$ and
$j^{-1}ij=i^{-1}$, while $\langle r,s\rangle\cong S_3$, since
$r^3=s^2=1$ and $srs=r^{-1}$. Their intersection is trivial and their
products exhaust the $48$ invertible matrices. The actual retained
$S_3$ complement acts faithfully on its normal $Q_8$, as is verified
by checking its centralizer on all eight elements. The explicit full
isomorphism above supplies the identification even without invoking a
classification of actions on $Q_8$.

\section{Designs and orientations}
For either choice of point type, the blocks are the neighborhoods of the
vertices of the other type. The certificate retains their complete point
sets, their constant size $k$, their number, and the number of distinct
neighborhoods. For each claimed strength $t$, the verifier counts, for every
$t$-subset $T$ of points, the blocks containing $T$. A constant value
$\lambda$ proves precisely the stated $S_\lambda(t,k,v)$ condition.
The first failed strength is certified by two explicit subsets with
unequal counts.

No larger nonvacuous strength can hold after such a failure. Indeed, for
a fixed $t$-subset $T$, count pairs $(x,B)$ with $x\notin T$ and
$T\cup\{x\}\subseteq B$. If every $(t+1)$-subset has multiplicity
$\lambda_{t+1}$, this gives
\[
 (k-t)\,\#\{B:T\subseteq B\}=(v-t)\lambda_{t+1},
\]
so the multiplicity at level $t$ is constant whenever $t<k$.

The point type in the design table is identified by its actual point count,
block size and multiplicity. The join records this orientation separately
from the ordering of the two subgroup names in Table 3. If the printed
invariants leave two rows indistinguishable, all permissible matches are
retained and a bijection is selected only within that ambiguity. No unique
embedding is inferred from an incomplete list of invariants.
Complete incidence is recorded separately: every block is the full point
set, and every subset is contained in all blocks. Whether such an object is
called a design depends on whether the convention requires proper or
distinct blocks. The literal counts make that convention visible.


\section{The source comparison and four orientation corrections}
The actual computation gives $309$ maximal subgroups and partitions all
$\binom{309}{2}=47586$ unordered pairs into $37$ orbits. The independent
checker validates $6278$ literal cosets and $28498$ subset counts in the
uniform design levels. Every asserted intersection name is joined through
its actual conjugator to an accepted structural proof. The extra matrix
certificate identifies the order-$48$ maximal subgroup.

Keeping the subgroup columns printed in Table 3, four triples require
correction:
\[
\begin{array}{c|cc}
\text{source row}&\text{printed}&\text{corrected}\\\hline
12,18 &[3,2,4]&[4,2,3]\\
23,35 &[3,3,4]&[4,3,3].
\end{array}
\]
Section 3 of the source defines $d_p$ as the eccentricity of a $G_1$ coset
and $d_l$ as that of a $G_2$ coset. The four printed triples instead use
the opposite orientation. This is a correction of the stated convention;
it does not change the abstract graphs. The original transcription and
its failed literal comparison are retained separately from the explicit
correction layer. All $70$ numerical design entries in Table 4 agree.

Row 35 has a short independent interpretation. Its point sets in the
opposite orientation have sizes $11$ and $55$, and the design is the full
set of two-element subsets of eleven points. Its incidence graph is the
subdivision of $K_{11}$. An original vertex has maximum distance three,
attained at a nonincident edge vertex. An edge vertex has maximum distance
four, attained at a disjoint edge vertex. The source puts the index-$55$
subgroup $M_9:2$ first and the index-$11$ subgroup $M_{10}$ second, so its
own convention gives $[4,3,3]$. Its gonality is three because a triangle
of $K_{11}$ becomes a cycle of length six. The same argument gives the
row 23 correction with $K_{12}$.

Rows 12 and 18 are likewise recovered with the stated type sizes. In the
orientation with $12$ or $11$ points, their point diameters are three and
block diameters four; the subgroup columns in Table 3 put the blocks first.
The full finite distance certificates establish both values. Table 4's
design orientation is retained independently of these subgroup columns.

Rows 10 and 11, and rows 16 and 17, have indistinguishable printed profiles.
The computation supplies two different pair orbits in each case. We retain
both possible source matches within each pair of rows, rather than pretend
that the printed invariants identify a unique embedding. This ambiguity
does not change any asserted table entry or the complete orbit census.

The omitted design rows 31 and 32 are respectively complete bipartite
graphs on parts of sizes $(12,55)$ and $(12,11)$. In each orientation every
block is the full point set. They satisfy the literal subset-counting
condition with repeated full blocks; their omission uses the usual
requirement of proper or distinct blocks, which is not stated in the
source's short definition. We record both the counts and this convention
qualification. No nontrivial design entry is lost.

\section{Attribution and boundaries}
The subgroup classification, its named structural results and the standard
properties of $M_{11}$ are borrowed with exact links to their accepted
components. The source supplies the assertions to be compared. Our work
here is the common finite certificate, its independent verification, the
explicit source joins and the matrix naming bridge. None of these finite
checks is presented as Lean verification or as a new compression measurement.
The complete source disposition requires the returned graph data and source
comparison to pass independent review.
\input{tables.tex}
\end{document}
